\documentclass[10pt,a4paper]{amsart}
\usepackage[margin=19mm]{geometry}
\usepackage{amsmath,amssymb,mathtools}
\usepackage[scr=rsfs,cal=euler]{mathalfa}
\usepackage{xfrac}
\usepackage[unicode,hidelinks,bookmarks=false]{hyperref}
\newcommand{\ie}{i.e.\ }
\newcommand{\cf}{cf.\ }
\newcommand{\resp}{respectively\ }
\newcommand{\Subsection}{\subsection}
\providecommand{\nobreakdash}{\nobreak\hbox{-}}
\setlength{\emergencystretch}{2em}
\multlinegap0pt
\usepackage{bm}
\usepackage{bbm}
\usepackage{dsfont}
\usepackage{xspace}
\usepackage{cancel}
\usepackage{mathtools}
\usepackage{amsthm}
\makeatletter
\@ifundefined{lemma}{\newtheorem{lemma}{Lemma}}{}
\makeatother
\def\E{{\mathcal E}}
\def\d{{\textnormal d}}
\def\r{{\mathbb R}}
\title[Mehler formula for Wronskians of Hermite polynomials]{Mehler formula for Wronskians of Hermite polynomials}
\author{Alexey Kuznetsov, Minjian Yuan}
\date{Source: arXiv:2606.02405v1 (CC BY 4.0)}
\begin{document}
\maketitle

\noindent\emph{Source and licence.} Mehler formula for Wronskians of Hermite polynomials. Version arXiv:2606.02405v1 (CC BY 4.0), distributed under the Creative Commons Attribution 4.0 International licence, as recorded in the arXiv licence metadata for this exact version and shown on the abstract page https://arxiv.org/abs/2606.02405v1. Full rights notice: licenses/arxiv-2606.02405-CC-BY-4.0.txt.\par\smallskip

\noindent\textit{Editorial scope.} Counted unit: the Lemma whose source label is \texttt{lemma6} in the article (source0.tex lines 685--692, source label \texttt{lemma6}), delivered together with the article's own complete original proof of it (source0.tex lines 693--802, one \texttt{proof} environment of 110 lines). No mathematics is written, repaired or re-derived: statement and proof are the article's own text line for line. Required context retained uncounted: def\_mathcal\_E (lines 306--309). Every definition, display and label of those lines is retained unchanged. Omitted from this package and not redistributed: 1--305, 310--684, 803--1235. Nothing inside a retained excerpt is omitted or altered: every retained body is delivered exactly as the article's lines are written, so no omission receipt of any kind accompanies this package. Citations: the retained excerpts carry no citation command at all, so no bibliography entry of the article is transcribed and no reference is redistributed. Reference adaptation: none was needed and none was made; every reference inside the delivered unit resolves inside it, so no unresolved reference or citation is produced. Printed slips kept literal: no printed word, symbol, hypothesis or number of the counted statement or of its proof was altered. Layout and class: the article's own class is \texttt{article} with packages including setspace, graphicx, amsmath, amsthm, amsfonts, amssymb, hyperref; the wrapper is the lane's pinned \texttt{amsart} editorial wrapper at 10pt on A4, which restates the theorem-like environments the retained text needs and adds the article's own macro definitions that the retained text needs (\texttt{\textbackslash E}, \texttt{\textbackslash d}, \texttt{\textbackslash r}), with extra wrapper packages (bm, bbm, dsfont, xspace, cancel, mathtools). The delivered numbering is the unit's own. Assets: no retained span contains a figure environment, a graphics-inclusion command, a PDF-page inclusion, a table or a TikZ picture; no image file is copied into this package, the package declares no asset dependency and no image file is redistributed with it. Licence: version arXiv:2606.02405v1 of this article is distributed under the Creative Commons Attribution 4.0 International licence, as recorded in the arXiv licence metadata for this exact version and shown on the abstract page https://arxiv.org/abs/2606.02405v1. A full rights notice is delivered at licenses/arxiv-2606.02405-CC-BY-4.0.txt. Characters: every retained excerpt is pure ASCII, so the delivered engine, XeLaTeX with its default Latin Modern text font, reports no missing character.
\begin{equation}\label{def_mathcal_E}
{\mathcal E}(z,x,y):=\sum\limits_{n\geqslant 0} z^n \frac{\psi_n(x) \psi_n(y)}{2^n n!}=
\frac{1}{\sqrt{1-z^2}} \exp\bigg( \frac{4xyz-(1+z^2)(x^2+y^2)}{2(1-z^2)} \bigg),
\end{equation}

\begin{lemma}\label{lemma6}
Let $m,n \in {\mathbb Z}_{\geqslant 0}$ and let $x, y$ be fixed real numbers satisfying $x>y$. Let $P(u)$ be a polynomial such that $P^{(i)}(y)=0$ for $i=0,1,\dots,n$. If $m\leqslant n+1$ then 
\begin{equation}\label{eqn_integral_limit}
\lim\limits_{z\to 1-} \frac{\partial_x^{m} \int_{-\infty}^y \E(z,x,u) e^{-u^2/2}  P(u) \d u}{\E(z,x,y)} 
=0. 
\end{equation}
If $m=n+2$, the above limit exists and is finite.  
\end{lemma}

\begin{proof}
First, we check that formula \eqref{def_mathcal_E} can be rewritten in the equivalent form  
\begin{equation}\label{E_factorization}
\E(z,x,y)  =e^{(x^2+y^2)/2  - 2xy/(1+z)} \times 
\zeta^{-1} 
e^{-(y-x)^2/\zeta^2 },
\end{equation}
where $\zeta=\zeta(z)=\sqrt{1-z^2}$. For the rest of this proof, we assume that $1/\sqrt{2}<z<1$, so that $0<\zeta<1/\sqrt{2}$.  


 
We claim that, for any $m\in {\mathbb Z}_{\geqslant 0}$ and any fixed $x,y \in \r$  such that $x>y$, we have
\begin{equation}\label{eqn_q_integral}
\zeta^{-2} e^{(y-x)^2/\zeta^2} \int_{-\infty}^y  
e^{-2xu/(1+z)-(u-x)^2/\zeta^2}   u^m \d u \to 
\frac{e^{-x y}y^m}{2(x-y)},  
\end{equation}
as $z\to 1-$ (note that $\zeta\to 0+$ as $z\to 1-$). 
Indeed, changing the variable of integration $u=y-\zeta^2 v$, we obtain
\[
\zeta^{-2}  e^{(y-x)^2/\zeta^2}  \int_{-\infty}^y  
e^{-2xu/(1+z)-(u-x)^2/\zeta^2}   u^m \d u
= e^{-2x y/(1+z)} \int_0^{\infty} e^{-2(x-y)v -\zeta^2 (v^2-2xv/(1+z))}  (y-\zeta^2 v)^m \d v,
\]
and the limit in \eqref{eqn_q_integral} follows by the dominated convergence theorem. 




Next, let $q(z,u)$ be a polynomial in $u$ of the form
\begin{equation}\label{form_of_q}
q(z,u)=\sum\limits_{i=0}^M a_i(z) u^i,
\end{equation}
where $M$ does not depend on $z$ and the coefficients $a_i(z)$ are continuous in some neighborhood of $z=1$. For $z\in (0,1)$ and $j \in {\mathbb Z}_{\geqslant 0}$ we consider 
\begin{equation}\label{def_Q_j}
{\mathcal Q}_{j}(z,x,y):=\zeta^{-1} \int_{-\infty}^y \Big[\partial_u^j  
e^{-(u-x)^2/\zeta^2} \Big] \times  e^{-2xu/(1+z)} q(z,u) (u-y)^{n+1} \d u. 
\end{equation}
From \eqref{E_factorization} and \eqref{eqn_q_integral} we conclude that 
\begin{equation*}
\frac{{\mathcal Q}_{0}(z,x,y)}{\E(z,x,y)} \to 0
\end{equation*}
as $z \to 1-$. When $0< j \leqslant n+1$, we 
integrate by parts $j$ times 
and obtain 
\begin{align}\label{Q_j_integration_by_parts}
{\mathcal Q}_j(z,x,y)&= 
\zeta^{-1} \sum\limits_{i=0}^{j-1} (-1)^i \Big[\partial_u^{j-1-i}  
e^{-(u-x)^2/\zeta^2} \Big]_{u=y} \times 
\Big[\partial_u^i \Big(  
e^{-2xu/(1+z)} q(z,u) (u-y)^{n+1}\Big)\Big]_{u=y} 
\\\nonumber &+(-1)^j \zeta^{-1} \int_{-\infty}^y 
e^{-(u-x)^2/\zeta^2} 
\partial_u^j  \Big(
e^{-2xu/(1+z)} q(z,u) (u-y)^{n+1}\Big) \d u.
\end{align}
In the finite sum in the above formula, we have $i\leqslant j-1 \leqslant n$, and therefore
\[
\Big[\partial_u^i  
e^{-2xu/(1+z)} q(z,u) (u-y)^{n+1}\Big]_{u=y}=0. 
\]
Thus, the finite sum in \eqref{Q_j_integration_by_parts} is equal to zero, and
it follows from \eqref{E_factorization} and \eqref{eqn_q_integral} that, for $0<j\leqslant n+1$,
\begin{equation*}
\frac{{\mathcal Q}_{j}(z,x,y)}{\E(z,x,y)} \to 0, 
\end{equation*}
as $z \to 1-$. 

When $j=n+2$, the preceding argument shows that  the terms with $0\leqslant i \leqslant j-2$ in the finite sum in \eqref{Q_j_integration_by_parts} are all equal to zero, and when $i=j-1=n+1$ we have
\[
\Big[\partial_u^{n+1} \Big(  
e^{-2xu/(1+z)} q(z,u) (u-y)^{n+1}\Big)\Big]_{u=y}= 
e^{-2xy/(1+z)} q(z,y) (n+1)!.
\]
Thus, it follows from 
\eqref{E_factorization}, \eqref{eqn_q_integral} and 
\eqref{Q_j_integration_by_parts} that, as $z\to 1-$,
\begin{equation*}
\frac{{\mathcal Q}_{n+2}(z,x,y)}{\E(z,x,y)} \to 
 (-1)^{n+1} 
 e^{-(x^2+y^2)/2} q(1,y) (n+1)!. 
\end{equation*}


Now we are ready to complete the proof of Lemma \ref{lemma6}. For any fixed $z\in (1/\sqrt{2},1)$, we interchange the integration and differentiation in the numerator in \eqref{eqn_integral_limit},   
use the factorization \eqref{E_factorization}, and obtain
\begin{equation}\label{numerator_as_sum_of_Q}
\partial_x^{m} \int_{-\infty}^y \E(z,x,u) e^{-u^2/2}  P(u) \d u
= \sum\limits_{j=0}^m (-1)^j \binom{m}{j} \widetilde {\mathcal Q}_{j}(z,x,y),
\end{equation}
where we have defined 
\begin{equation}\label{def_tilde_Q}
\widetilde {\mathcal Q}_{j}(z,x,y):=\zeta^{-1} \int_{-\infty}^y
\Big[ \partial_u^{j} e^{-(u-x)^2/\zeta^2} \Big]  \times \Big[ \partial_x^{m-j} e^{x^2/2-2xu/(1+z)} \Big]  \times P(u) \d u.
\end{equation}
When deriving the above formula, we used the fact that 
\[
\partial_x^{j} e^{-(u-x)^2/\zeta^2}=(-1)^j \partial_u^{j} e^{-(u-x)^2/\zeta^2}. 
\]
The interchange of integration and differentiation is justified since, for $z\in (1/\sqrt{2},1)$, we have $1/\zeta^2>2$; thus, for all $u\in \r$ the integrand in \eqref{def_tilde_Q} is dominated by $C\exp(- u^2)$ for some $C>0$, uniformly in $x$ on compact subsets of $\r$.  We note that the conditions  $P^{(i)}(y)=0$ for $i=0,1,\dots,n$ imply that $P(u)=q_1(u) (u-y)^{n+1}$ for some polynomial $q_1$, which in turn implies
\[
\Big[ \partial_x^{m-j} e^{x^2/2-2xu/(1+z)} \Big] \times  P(u)=
e^{x^2/2-2xu/(1+z)} q_2(x,u/(1+z),u) (u-y)^{n+1},
\]
where $q_2(x_1,x_2,x_3)$ is some polynomial in three variables.   The important conclusion is that $q_3(z,u)=q_2(x,u/(1+z),u)$ is of the form \eqref{form_of_q} (recall that everywhere in this proof $x$ is fixed, so dependence of the coefficients of the polynomial $q_3(z,u)$ on $x$ is not a concern). Thus, $\widetilde {\mathcal Q}_{j}(z,x,y)$ is an integral of the form \eqref{def_Q_j}, and using the results about ${\mathcal Q}_j(z,x,y)$ that we established above, we conclude that the limit  
\[
\lim\limits_{z\to 1-} \frac{\widetilde {\mathcal Q}_{j}(z,x,y)}{\E(z,x,y)} 
\]
exists and is equal to zero for all $0\leqslant j \leqslant n+1$, while for $j=n+2$ the limit exists and is finite. The desired result now follows from \eqref{numerator_as_sum_of_Q}.
\end{proof}

\end{document}
